\chapter{Implementation and Validation}
\section*{Introduction}
This chapter presents the implementation of the principal workflows and the
technical environment used during the internship. It connects the design
decisions to concrete backend services, Angular screens, persistence rules and
API operations.

\section{Development environment and technologies}
\begin{table}[H]\centering
\caption{Technologies and versions}\label{tab:technologies}
\begin{tabularx}{\textwidth}{>{\bfseries}p{0.25\textwidth}p{0.20\textwidth}X}
\toprule Technology & Version & Purpose\\\midrule
Angular & 21 & Single-page frontend application.\\
TypeScript & 5.9 & Frontend implementation language.\\
Spring Boot & 3.3.5 & REST backend framework.\\
Java & 21 target & Backend implementation language.\\
Spring Data JPA & Project version & Relational persistence and repositories.\\
MySQL & Project environment & Main relational database.\\
H2 & Test environment & Isolated backend test database.\\
Maven & Project environment & Backend dependency and build management.\\
Node.js / npm & Project environment & Frontend dependency and build management.\\
Gemini API & Gemini Flash & AI-assisted quiz generation.\\
OpenAPI / Swagger & Project dependency & Interactive API documentation and testing.\\
\bottomrule
\end{tabularx}
\end{table}

\section{Backend implementation}
The backend is implemented with Spring Boot 3.3, Spring Web, Spring Data JPA,
validation and Spring Security. Controllers expose focused REST resources for
authentication, courses, materials, quizzes, progress and administration.
Services contain the business rules and repositories isolate database access.



\begin{table}[H]\centering
\caption{Main backend layers}\label{tab:backend-layers}
\begin{tabularx}{\textwidth}{>{\bfseries}p{0.22\textwidth}X}
\toprule Layer & Responsibility\\\midrule
Controller & HTTP mapping, request validation and response status.\\
Service & Lifecycle rules, ownership, grading and transactional operations.\\
Repository & JPA queries and dependent-record lookup.\\
Entity & Relational domain model and associations.\\
Integration & Gemini client, document extraction and file serving.\\\\
\bottomrule
\end{tabularx}
\end{table}

\section{Frontend implementation}
The frontend is an Angular 21 single-page application styled with Tailwind
utilities and a shared Safyron design system. Shared navigation, cards,
status badges, tables and action controls are reused across role dashboards.
The interface deliberately exposes different actions to learners, instructors
and administrators.

\section{Learner experience}
\img{assets/screenshots/login.png}{Live screenshot captured from the running login page}
\img{assets/screenshots/learner-dashboard.png}{Live screenshot captured from the learner portal}
The learner portal presents published modules with their completion state,
material controls and quiz availability. Files can be opened in the browser or
downloaded. Once a quiz has been submitted, its action is visually disabled
while the learning material remains available.

\section{Instructor studio}
\img{assets/screenshots/instructor-dashboard.png}{Live screenshot captured from the instructor studio}
The instructor studio combines module management, document intake, AI quiz
synthesis and manual authoring. It also presents pending submissions with
learner and module context. The detailed view expands the selected submission
to show each question, the learner answer and the correct answer.

\section{Administrator console}
\img{assets/screenshots/admin-dashboard.png}{Live screenshot captured from the administrator console}
The administrator console provides user creation and deletion, module
oversight, quiz management and submission indicators. The grouped view makes
the module the primary reporting unit and includes the submission date and
review state.

\section{Answer persistence and grading}
Answers are stored using the persistent question identifier rather than the
position of a question in the displayed array. This distinction prevents
review screens from displaying ``Not answered'' when the learner selected
option zero or when questions were reordered. The grading service stores the
submitted answer map as JSON and creates one result per learner and quiz.

\section{File management and deletion}
The material controller exposes an inline file endpoint and a download action.
Course deletion is transactional and removes quiz results, progress, quizzes
and materials before deleting the module. User deletion similarly removes
learner-dependent records and detaches owned entities. These measures prevent
the foreign-key error that originally appeared when deleting users directly
from the database.

\section{Swagger and API usage}
OpenAPI documentation is exposed by the backend and provides an interactive
Swagger UI for testing. Typical operations include:
\begin{lstlisting}
POST /api/auth/login
GET  /api/courses/published
POST /api/quizzes/generate
POST /api/quizzes/manual
POST /api/progress/quiz/{quizId}/submit
PUT  /api/quizzes/results/{resultId}/approve
PUT  /api/quizzes/results/{resultId}/unapprove
\end{lstlisting}

\section{Implementation difficulties and corrective actions}
Several defects were discovered through integration and browser validation.
Learner answers initially used array positions while the instructor review
expected persistent question identifiers; the answer binding was corrected and
legacy values were normalized. A second quiz attempt was initially possible,
so the service now rejects an existing learner--quiz result with HTTP 409.

Module and user deletion initially failed because dependent rows remained in
the database. Transactional cleanup, repository deletion methods and
database-level cascade or ``SET NULL'' behavior were added. AI generation also
produced misleading forbidden responses when an application exception was
forwarded to a protected error route; explicit error responses now expose the
actual failure to the client.

\section{Validation and Results}
\section*{Validation introduction}
Validation was performed at several levels because a successful compilation
does not prove that the learning lifecycle is correct. The checks below cover
automated backend tests, frontend builds, API behavior, database consistency
and browser journeys for the three roles.

\section{Validation strategy}
Validation combined unit and integration tests for the backend, an Angular
production build, Swagger checks and live HTTP scenarios. The objective was
not only to verify that an endpoint returned a response, but to verify the
business rule represented by that response.

\begin{table}[H]\centering
\caption{Validation activities}\label{tab:validation}
\begin{tabularx}{\textwidth}{>{\bfseries}p{0.30\textwidth}X}
\toprule Activity & Result\\\midrule
Backend clean tests & Passed after CRUD, grading and deletion changes.\\
Angular production build & Passed with the final dashboard and branding changes.\\
OpenAPI availability & Swagger UI and OpenAPI JSON returned successfully.\\
Module deletion & Returned successful deletion after dependent cleanup.\\
Answer persistence & Review returned the selected option for the correct question.\\
Duplicate submission & A second attempt was rejected with HTTP 409.\\
User deletion & Dependent learner records were removed without a foreign-key failure.\\
\bottomrule
\end{tabularx}
\end{table}

\section{Validation traceability}
\begin{table}[H]\centering
\caption{Traceability between requirements and validation}
\label{tab:traceability}
\begin{tabularx}{\textwidth}{>{\bfseries}p{0.25\textwidth}X}
\toprule Requirement & Validation evidence\\\midrule
Security & JWT-aware login design, role-aware routes and protected operations.\\
Consistency & Real module and user deletion scenarios without foreign-key failure.\\
Traceability & Detailed answer review, submission date and approval/unapproval.\\
Maintainability & Successful layered backend tests and Angular production build.\\
AI continuity & Manual quiz creation remains available beside AI generation.\\
One-time assessment & First submission succeeds and repeated submission returns 409.\\
Material availability & Files remain accessible after quiz submission.\\
\bottomrule
\end{tabularx}
\end{table}

\section{Functional scenarios}
The first scenario creates a module and attaches documents. The second
generates or manually creates a quiz and publishes it. The learner then opens
the module, selects answers and submits once. The instructor sees the
submission in the review area, expands the detailed answers and approves it.
The administrator can consult the grouped submission statistics.

\section{Quality results}
The principal defects discovered during development were incorrect answer
keys, misleading HTTP 403 responses caused by error forwarding, incomplete
dependent deletion and an interface that allowed a second attempt. Each issue
was fixed at its root: question IDs are now used for answer maps, application
errors return explicit statuses, deletion is transactional, and both backend
and frontend prevent retakes.

\section{Security and operational limitations}
The JWT integration provides the intended authentication and authorization
boundary for deployment. Secret values, including the AI provider key, must
remain in environment configuration and must be rotated if exposed. File
uploads should receive size, type and content scanning limits before production.
Rate limiting and audit retention are also recommended for a public
deployment.

\section{Difficulties, limitations and perspectives}
The most important difficulties were not isolated coding errors but
interactions between features. Deleting a module required understanding all
dependent records. Persisting an answer required a stable question key rather
than a temporary array index. The approval workflow required a reversible
state, not only a one-way button. These issues were corrected through service
rules, repository operations and end-to-end checks.

The current validation does not constitute a production load test or an
independent security audit. Future work should measure performance with a
larger catalogue, add automated accessibility checks, scan uploaded files,
rotate refresh tokens, export audit records and add notifications. The AI
workflow should also display confidence or source references and be evaluated
against a larger curated question corpus.

\section{User experience observations}
The live screenshots show a coherent visual language: blue-gray surfaces,
strong contrast, shield branding, status badges and action-oriented cards.
The instructor screen emphasizes operational work, while the learner screen
keeps the path to course material and assessment simple. The administrator
screen prioritizes counts, grouped submissions and governance actions.

\section{Personal contribution}
This internship required work across requirements analysis, Angular
components, Spring services, relational integrity, third-party AI integration,
API documentation and testing. The most valuable lesson was that a feature is
not complete when its happy path works: its deletion behavior, error response,
role boundary and persistence format must also be verified.




\section*{Conclusion}
The implementation connects the role-specific Angular experience to a
transactional Spring Boot backend, and the validation activities confirm the
main learner, instructor and administrator workflows. The report also
identifies the remaining production-hardening work without overstating the
evidence obtained during a one-month internship.
